\subsection{Large Language Models}\label{subsec:large-language-models}
Generalized language model was created from the demand for tasks with increasing complexity in natural language.
These tasks may include translation, summarization, information retrieval, code generation, and many more~\parencite[2]{naveed2025llm}.
Large Language Models are only subset of the field called Artificial Intelligence specialized in natural language processing.

The first models in natural language processing used rule-based techniques dominant in the 1950s and statistical models
in the 1980s.
The idea of Artificial Neural Network existed in the 1940s but hasn't seen much development.
These models were mainly used for binary classification, but they lacked the understanding of textual context
to deal with complex tasks~\parencite[History of Large Language Models]{raiaan2024llm}.

Large Language Models can be dissected into components each providing crucial functionality
~\parencite[Background of Large Language Models]{raiaan2024llm}:

\begin{itemize}
    \item Tokenization - The initial text is parsed into a sequence of tokens such as characters, symbols, parts of words, full words, parts of sentences.
        The main algorithms used in tokenization include WordPiece, UnigramLM, and Byte Pair Encoding.
    \item Attention - This mechanism improves performance of the models and is responsible for figuring out representation
        of input sequences and linking between tokens in the sequence.
    \item Activation Function - The function gives Large Language Model the ability to learn non-trivial concepts.
        Some of the activation functions are ReLU, GeLU, and other GLU variations.
    \item Normalization layer - Essential part in Large Language Model's convergence rate and improves stability during training.
        Some of the approaches are DeepNorm, RMSNorm.
    \item Data processing - Important part of the training process is general data processing.
        This ensures that the data is filtered for quality, de-duplicated, and removes personal data,
        which could lead to privacy concerns.
\end{itemize}

Usage of Large Language Models for generative or transformative purposes leads to ethical problems.
The organization should be transparent about the extent of AI usage in their works
to ensure expectation about quality of the provided service~\parencite[Ethical and Legal Considerations in AI Translation]{ablaikhan_translation_ai}.

\subsubsection{Translation}
Recent advances in Large Language Models have significantly improved the quality of translations.
Unlike traditional statistical or rule-based translation systems that are built on task-specific architectures,
Large Language Models are trained on large-scale multilingual dataset gives the ability to understand the textual context
and translate it accordingly to the meaning.

Classic GPT-based models excel at producing naturally sounding translations and are not limited to a few phrases,
which gives them ability to translate entire documents while preserving context.
Where human post-editing is still needed is for idiomatic phrases, ambiguous phrases, or translating words with multiple meanings.
Since Large Language Models are probabilistic their output can be different for the same sample of text when one output is strongly desired.
This probabilistic nature also introduces biases from training data,
particularly found in gendered languages and when translating sensitive topics~\parencite[Quality of AI Translations: Successes and Gaps]{ablaikhan_translation_ai}.
Models can be fine-tuned in prompt to translate the text in certain ways,
these translations are closer to transformation of a text and human intervention is needed to check for concept preservation.

The use of Large Language Models for translation introduces ethical and legal concerns.
Translations are derivative works, meaning that copyright of the original content must be respected,
even when Large Language Models is used for translation.
The user or organization using the translation tool is responsible for correcting translation mistakes introduced by the tool.

Current systems are efficient at dealing with routine translation tasks, such as translating phrases without context,
basic text, or structured text.
However, the same system still need human surveillance to meet the quality standards of translations~\parencite[AI vs. Human Translators: Can AI Replace Them?]{ablaikhan_translation_ai}.

\subsubsection{Code Generation}
Recent advances in Large Language Models have significantly improved the quality of code generations.
Code generation tasks are defined as generating code that performs the actions described in a given specification.
The description can be written as natural language prompt or introduce domain specific pseudo-language.

Large Language Models generate code by analyzing vast code bases in their training dataset~\parencite[AI Models and Techniques for Code Generation]{khade2025review_of_llm_code}.

Some of the common tasks may include~\parencite[Introduction]{chen2024survey_codellm}:

\begin{itemize}
    \item Prompt to executable code - The process of converting a natural language request into executable code.
        The generated code can be of different types: a small procedure with a well-defined functionality (e.g., a sorting algorithm),
        an architectural component implementing logic for a specific domain (e.g., a custom collection),
        or an executable script for an application or an operating system shell.
    \item Interpolating code templates - The prompt consists of mostly finished code snippet and description how to fill
        in the holes from context.
        Potential usage may be found in translation of UI templates.
    \item Refactoring existing code - The input includes the source code to be refactored and additional context.
        The context may include more source code or prompt describing the architectural ideas.
        The input also includes how the given code should be transformed.
        It is expected that the output is executable code, which interface and functionality is kept the same as the original code.
    \item Test case generation - The input includes the source code to be tested, additional context, and testing framework/library.
        The output are test cases validating the performance of the code in edge case situations and correctness in expected situations.
\end{itemize}

The usage of AI for code generation enables the developers to offload repetitive tasks, and they can focus on architectural design,
critical optimization problems, and other demanding tasks.
The shift from pure coding jobs to AI-assisted software engineering has reshaped the job market.
This development raises concerns about job displacement and changing demands for software developers~\parencite[Changes in Software Engeneering Job Roles]{khade2025review_of_llm_code}.

\subsubsection{Code Documentation}

